\section{Related Work}
\label{sec:related}

A starting point is lexical retrieval where probabilistic ranking models such as BM25 remain strong baselines thanks to their efficiency, interpretability and robustness in candidate generation \citep{RobertsonZaragoza2009}.

Neural retrieval methods have been proposed to mitigate the lexical mismatch between queries and documents by learning dense semantic representations. More broadly, neural information retrieval has shifted attention from exact token overlap toward representation learning and learned matching functions \citep{MarchesinEtAl2020}. Within this line of work, dense dual-encoder architectures such as Dense Passage Retrieval (DPR) demonstrated that learned query and document embeddings can be highly effective for first-stage retrieval \citep{KarpukhinEtAl2020}. This is especially relevant to our task because social media posts often paraphrase scientific findings rather than reuse the exact wording found in the original publication. Building on this direction, BGE-M3 unified dense, sparse, and multi-vector retrieval in a single multilingual model \citep{bge-m3},
enabling hybrid scoring that combines embedding similarity with lexical token weights, a combination particularly well-suited to cross-lingual settings where queries and documents share few surface tokens.

Multilingual information retrieval presents additional challenges beyond monolingual settings, as systems must bridge language mismatches before relevance can be estimated. Machine translation provides one practical solution, but it can introduce errors and vocabulary drift, especially for technical or domain-specific terminology \citep{SennrichEtAl2016}.

Vocabulary mismatch between short, informal queries and longer scientific documents can be further reduced through query expansion. Classical pseudo-relevance feedback methods such as RM3 \citep{LavrenkoEtAl2001} address this issue by re-weighting query terms using top-ranked documents from an initial retrieval pass. Neural extensions of this paradigm replace or complement lexical evidence with dense representations \citep{WangEtAl2021PRF}. Following this general approach, our pipeline incorporates embedding-based query expansion to enrich query representations before scoring against the corpus.

Between initial retrieval and final ranking, many modern systems adopt multi-stage pipelines to balance efficiency and effectiveness \citep{MarchesinEtAl2020}. Within these pipelines, late-interaction models such as ColBERT preserve fine-grained token-level matching while remaining substantially more scalable than fully pairwise neural scoring over the whole collection \citep{khattab2020colbert}. Similarly, transformer-based re-rankers have become a common second-stage component for refining an initial candidate set through stronger query-document interactions \citep{NogueiraEtAl2020}. Taken together, these approaches motivate retrieval architectures that first generate a manageable set of candidates and then apply more expensive relevance estimation only where it is most useful.

Recent work has argued that reranker fine-tuning benefits from hard negatives because they force the model to learn finer-grained relevance
boundaries than easy or random negatives provide \citep{MeghwaniEtAl2025HardNegatives}. Adapting large language models to retrieval tasks is further facilitated by parameter-efficient fine-tuning with LoRA \citep{hu2022lora}, which has been shown to match or approach full fine-tuning quality for passage reranking while substantially reducing GPU memory requirements \citep{nijasure2025relevance}. This perspective is also consistent with recent practical guidance for training cross-encoder rerankers with Sentence Transformers \citep{Aarsen2025TrainReranker}.
